In this section we will use the theory built up in Sections \ref{sec:Goodwillie} and \ref{sec:ThomPushout} to reduce Theorem \ref{thm:intro-main} to a concrete Toda bracket computation.
The final result of this section, \Cref{lem:toda-main}, is the only statement from Sections \ref{sec:Goodwillie}-\ref{sec:Toda} that is used later in the paper. The lemma expresses the Toda bracket in an explicit enough form that we will be able to bound its $\HFp$-Adams filtrations in \Cref{sec:SyntheticToda}.
%\Cref{lem:toda-main} gives an explicit form to the Toda bracket which will make it possible for us to bound its $\HFp$-Adams filtration in \Cref{sec:SyntheticToda}.

Recall once more the fundamental square
$$
\begin{tikzcd}[column sep = huge]
\sOf^{\otimes 2} \arrow{dd}{m} \arrow{r}{1 \otimes J} & \sOf \arrow{dd}{J} \\ \\
\sOf \arrow{r}{J} \arrow[Leftrightarrow]{ruu}{a} & \bS.
\end{tikzcd}
$$
Let $P$ denote the  cofiber
$$
\begin{tikzcd} [column sep = huge]
P=\mathrm{cofiber}(\sOf^{\otimes 2} \arrow{r}{1 \otimes J - m} & \sOf).
\end{tikzcd}
$$
Then, as explained in Remark \ref{cnstr:square}, the homotopy $a$ gives rise to a canonical map $P \to \mathbb{S}$.  According to Theorem \ref{thm:cofiber}, there is a long exact sequence
\begin{equation} \label{eq:les}
  \pi_{8n-1} (P) \to \pi_{8n-1} (\mathbb{S}) \to \pi_{8n-1} \left(\MO \langle 4n \rangle\right) \to \pi_{8n-2} (P) \to \pi_{8n-2} (\mathbb{S}),
\end{equation}
which we will use to compute $\pi_{8n-1} \left( \MO \langle 4n \rangle \right)$.

\begin{lem} \label{lem:unit-surjectivity}
    The group $\pi_{8n-2} (P)$ is trivial.
\end{lem}

\begin{proof}
Consider the long exact sequence
\begin{center}
\begin{tikzcd}[column sep=tiny]
		\pi_{8n-2} \left( \sOf^{\otimes 2} \right) \arrow[rr] & \arrow[d, phantom, ""{coordinate, name=Z}] & \pi_{8n-2} \left( \sOf \right) \arrow[dll, rounded corners, to path={ -- ([xshift=2ex]\tikztostart.east) |- (Z) [near end]\tikztonodes -| ([xshift=-2ex]\tikztotarget.west) --(\tikztotarget)}] \\
%, to path={ -- ([xshift=2ex]\tikzctostart.east) |- (Z) [near end]\tikztonodes -| ([xshift=-2ex]\tikztotarget.west) --(\tikztotarget)}
		\pi_{8n-2} (P) \arrow[rr] & \vphantom{a} & \pi_{8n-3} \left( \sOf^{\otimes 2} \right) \cong 0.
	\end{tikzcd}
\end{center}
%$$\pi_{8n-2} \left( \sOf^{\otimes 2} \right) \to \pi_{8n-2} \left( \sOf \right) \to \pi_{8n-2} (P) \to \pi_{8n-3} \left( \sOf^{\otimes 2} \right) \cong 0.$$
According to Corollary \ref{cor:gen-by-x2}, $\pi_{8n-2} \left( \sOf \right) \cong \mathbb{Z}/2\mathbb{Z}$, generated by $x^2$.
We will thus be done upon showing that $x^2$ is in the image of the map
$$\pi_{8n-2} \left( \sOf^{\otimes 2} \right) \to \pi_{8n-2} \left( \sOf \right).$$
The class $x \otimes x$ in the domain is sent to $xJ(x)-x^2$.  By Lemma \ref{lem:xJ(x)null}, $xJ(x)=0$.
\end{proof}

To complete the proof of Theorem \ref{thm:intro-main}, it remains to compute the image of the map
$$\pi_{8n-1} (P) \to \pi_{8n-1} (\mathbb{S}).$$
Note that the definition of $P$ as the cofiber of a map $\sOf^{\otimes 2} \to \sOf$ means that there is a canonical map
$\pi_{8n-1} (P) \to \pi_{8n-2} \left( \sOf^{\otimes 2} \right).$

\begin{lem} \label{lem:Reduction1}
Suppose $\ell$ is any class in $\pi_{8n-1} P$ which maps to 
$$2(x \otimes x) \in \pi_{8n-2} \left( \sOf^{\otimes 2} \right) \cong \Z\{x \otimes x\}.$$
Then $\J_{8n-1}$ and the image of $\ell$ in $\pi_{8n-1} (\mathbb{S})$ generate the kernel of the map
$$\pi_{8n-1} (\mathbb{S}) \to \pi_{8n-1} (\MO \langle 4n \rangle).$$
%If the image of $\ell$ in $\pi_{8n-1} (\mathbb{S})$ is an element of $\J_{8n-1}$, Theorem \ref{thm:intro-main} will follow.
\end{lem}

\begin{proof}
    %To prove \Cref{thm:intro-main}, we must show that, under the map $P \to \Ss$, every element of $\pi_{8n-1}P$ lands in $\J_{8n-1}$.
  By equation (\ref{eq:les}), it suffices to show that the image of $\pi_{8n-1} (P) \to \pi_{8n-1} (\mathbb{S})$ is generated by $\J_{8n-1}$ and $\ell$.
  We will argue using the cofiber sequence
  $$ \sOf \to P \to \Sigma \sOf^{\otimes 2}. $$
  The composite map $\sOf \to P \to \mathbb{S}$ is by definition $J$. Therefore, by Theorem \ref{thm:imJwelldefined}, it has image exactly $\mathcal{J}_{8n-1}$ in degree $8n-1$. What remains is to show that $2 (x \otimes x)$ generates the subgroup of elements of $\pi_{8n-2} \left(\sOf^{\otimes 2} \right)$ that lift to $\pi_{8n-1} (P)$. Consider the map
  $$ \Z\{x \otimes x\} \cong \pi_{8n-2} \left(\sOf^{\otimes 2}\right) \to \pi_{8n-2} \left(\sOf\right) \cong (\Z/2)\{x^2\} $$
  The class $x \otimes x$ maps to 
    $$x^2 -x J(x) = x^2 \in \pi_{8n-2} \left( \sOf \right),$$
    since $x J(x) = 0$ by \Cref{lem:xJ(x)null}. Therefore, $2(x \otimes x)$ is a generator of the subgroup of elements which lift.
\end{proof}

Unwinding the definition of $P$, it is helpful to restate Lemma \ref{lem:Reduction1} in the following equivalent form:

\begin{cnstr} \label{cnstr:toda-basic}
Recall from \Cref{cor:gen-by-x2} and \Cref{lem:xJ(x)null} that $2xJ(x)=2x^2=0$ in $\pi_{8n-2} \left(\spOf\right)$.  We may therefore choose (completely arbitrary) nullhomotopies $f$ and $b$ completing the following diagram: 
$$
\begin{tikzcd}[column sep = huge]
\bS^{8n-2} \arrow{r}{2(x \otimes x)}
\arrow[rdd, bend right=20,""{name=D}] \arrow[rdd, bend right=20,swap,"0"] \arrow[rr,bend left = 30,""{name=U},"0"]
& \sOf^{\otimes 2} \arrow[Leftrightarrow, from=D, "b"] \arrow[Leftrightarrow, from=U, "f"] \arrow{dd}{m} \arrow{r}{1 \otimes J}
& \sOf \arrow{dd}{J}  \\ \\
& \sOf \arrow[Leftrightarrow]{ruu}{a} \arrow{r}{J}  & \bS. 
\end{tikzcd}
$$
Composing all three of these homotopies yields a homotopy between the map
$$0:\bS^{8n-2} \to \bS$$
  and itself, or equivalently a loop in the pointed mapping space $\mathrm{Hom}_{*}(\bS^{8n-2},\bS)$, or an element $z \in \pi_{8n-1} \bS$.  This Toda bracket $z$ is well-defined up to changing the nullhomotopies $f$ and $b$. The sets of homotopy classes of nullhomotopies $f$ and $b$ are torsors for $\pi_{8n-1} \sOf$, and changing either $f$ or $b$ by an element $y \in \pi_{8n-1} \sOf$ has the effect of changing the element $z$ by $J(y)$. The class $z$ therefore has indeterminancy equal to $\pi_{8n-1}J$, which is equal to $\J_{8n-1}$ by \Cref{thm:imJwelldefined}.
\end{cnstr}

\begin{lem} \label{lem:toda-basic}
The kernel of the map
$$\pi_{8n-1} (\mathbb{S}) \to \pi_{8n-1} (\MO \langle 4n \rangle)$$
is generated by $\J_{8n-1}$ and the Toda bracket $z$ of \Cref{cnstr:toda-basic}.
\end{lem}

\begin{proof}
The nullhomotopies $f$ and $b$ combine to give a nullhomotopy of the composite
$$
\begin{tikzcd}[column sep = huge]
\bS^{8n-2} \arrow{r}{2(x \otimes x)} & \sOf^{\otimes 2} \arrow{r}{1 \otimes J -m} & \sOf,
\end{tikzcd}
$$
which is exactly the data of a lift of $2(x \otimes x)$ to a class in
\[
\begin{tikzcd} [column sep = huge]
    \pi_{8n-2}\mathrm{fib}(\sOf^{\otimes 2} \arrow{r}{1 \otimes J - m} & \sOf) = \pi_{8n-2} \left( \Sigma^{-1} P \right) \cong \pi_{8n-1} (P).
\end{tikzcd}
\]
The conclusion therefore follows from \Cref{lem:Reduction1}.
\end{proof}

Our strategy will be to choose the nullhomotopies $f$ and $b$, or equivalently the lift $\ell$ in Lemma \ref{lem:Reduction1}, as judiciously as possible. It will be because of these choices that we will be able to establish our $\HFp$-Adams filtration bounds in \Cref{sec:SyntheticToda}. Let us begin by making a careful choice of the nullhomotopy $b$:

\begin{rec}
Suppose that $R$ is a homotopy commutative ring spectrum, and $r$ an element of $\pi_{2*+1} R$.  Then the graded commutativity of $\pi_*(R)$ ensures that $2r^2 =0$ in $\pi_{4*+2}(R)$.
  A small part of the data of an $\E_\infty$-structure on $R$ is a \emph{canonical} nullhomotopy of $2r^2$.
  Indeed, given $r : \bS^{2n+1} \to R$, there is a canonical factorization of $r^2$ as
  \[\bS^{4n+2} \to D_2 (\bS^{2n+1}) \to D_2 (R) \to R.\]
  The canonical nullhomotopy of $2r^2$ arises from fixing an identification of the $(4n+3)$-skeleton of $D_2 (\bS^{2n+1})$ with $\Sigma^{4n+2} C(2)$.
\end{rec}

\begin{cnstr} \label{cnstr:toda-main-preskeleton}
Let $h$ denote the canonical nullhomotopy of $2J(x)^2$ that arises from the fact that $J(x) \in \pi_{4n-1} \bS$ is an element in the odd degree homotopy of the $\E_{\infty}$-ring $\bS$.  Let $g$ denote the canonical homotopy $J(xJ(x))\simeq J(x)^2$ that arises from $J$ being a map of right $\bS$-modules, and let $f$ denote a completely arbitrary nullhomotopy of $2xJ(x)$.  Then we may form the following diagram,
$$
\begin{tikzcd}[column sep = huge]
\mathbb{S}^{8n-2} \arrow{r}{2}
\arrow[rrdd, bend right=20,""{name=D}] \arrow[rrdd, bend right=20,swap,"0"] \arrow[rr,bend left = 30,""{name=U},"0"]
\arrow[rrdd, bend right=20,""{name=D}] \arrow[rrdd, bend right=20,swap,"0"] \arrow[rr,bend left = 30,""{name=U},"0"]
& \mathbb{S}^{8n-2} \arrow[Leftrightarrow, from=D, "h"] \arrow[Leftrightarrow, from=U, "f"] \arrow{r}{xJ(x)} \arrow[rdd,""{name=MU},"J(x)^2" description] \arrow[rdd,swap, ""{name=MD},"J(x)^2" description]
& \sOf \arrow{dd}{J}  \arrow[bend left=20,Leftrightarrow, from=MU, swap, "g"] \\ \\
& & \mathbb{S}, 
\end{tikzcd}
$$
which composes to give a Toda bracket $w \in \pi_{8n-1} \bS$.
\end{cnstr}

\begin{lem} \label{lem:toda-main-preskeleton}
Let $w$ denote the Toda bracket of Construction \ref{cnstr:toda-main-preskeleton}.
Then the kernel of the map
$$\pi_{8n-1} (\mathbb{S}) \to \pi_{8n-1} (\MO \langle 4n \rangle)$$
is generated by $\J_{8n-1}$ and $w$. 
\end{lem}

\begin{proof}
Composing (whiskering) the homotopy $a$
$$
\begin{tikzcd}[column sep = huge]
\sOf^{\otimes 2} \arrow{dd}{m} \arrow{r}{1 \otimes J} & \sOf \arrow{dd}{J} \\ \\
\sOf \arrow{r}{J} \arrow[Leftrightarrow]{ruu}{a} & \bS.
\end{tikzcd}
$$
along the map $\mathbb{S}^{8n-2} \stackrel{x \otimes x}{\longrightarrow} \sOf^{\otimes 2}$ yields the diagram
$$
\begin{tikzcd}[column sep = huge]
\bS^{8n-2} \arrow{dd}{x^2} \arrow{r}{xJ(x)} \arrow[rdd,""{name=MU},"J(x)^2" description] \arrow[rdd,swap, ""{name=MD},"J(x)^2" description] &
\sOf \arrow{dd}{J} \arrow[bend left=10,Leftrightarrow, from=MU, swap, "g"] \\ \\
\sOf \arrow{r}{J} \arrow[bend left=10,Leftrightarrow, from=MD, swap, "c"]  & \bS. 
\end{tikzcd}
$$
Here, $g$ is the homotopy from Construction \ref{cnstr:toda-main-preskeleton}, and $c$ is the natural homotopy arising from the structure of $J$ as a ring homomorphism.  Consider now the slightly extended diagram 
$$
\begin{tikzcd}[column sep = huge]
\bS^{8n-2} \arrow{r}{2} &\bS^{8n-2} \arrow{dd}{x^2} \arrow{r}{xJ(x)} \arrow[rdd,""{name=MU},"J(x)^2" description] \arrow[rdd,swap, ""{name=MD},"J(x)^2" description] &
\sOf \arrow{dd}{J} \arrow[bend left=10,Leftrightarrow, from=MU, swap, "g"] \\ \\
&\sOf \arrow{r}{J} \arrow[bend left=10,Leftrightarrow, from=MD, swap, "c"]  & \bS. 
\end{tikzcd}
$$
To put ourselves in the situation of Lemma \ref{lem:toda-basic}, we must choose a nullhomotopy $f$ of $2xJ(x)$ as well as a nullhomotopy $b$ of $2x^2$.  The result follows from choosing $b$ to the canonical nullhomotopy of $2x^2$ arising from the fact that $\sOf$ is a (non-unital) $\mathbb{E}_\infty$-ring spectrum.  Since $J$ is naturally a map of $\mathbb{E}_\infty$-rings, and not just of $\mathbb{A}_\infty$-rings, this canonical nullhomotopy of $2x^2$ will compose with $c$ to be the canonical nullhomotopy $h$ of $2J(x)^2$.
\end{proof}

We record one final technical reduction to end this section.

\begin{dfn} \label{dfn:skeleton}
Let
	$$M \longrightarrow \Sigma^{\infty} \Oo \mathrm{\langle 4n-1 \rangle}$$
	denote the inclusion of an $(8n-1)$-skeleton of $\Sigma^{\infty} \Oo \mathrm{\langle 4n-1 \rangle}$. By the inclusion of an $(8n-1)$-skeleton, we mean in particular that the induced map
	$$\mathrm{H}_*(M;\F_p) \longrightarrow \mathrm{H}_*(\Sigma^{\infty} \Oo \mathrm{\langle 4n-1 \rangle};\F_p)$$
is an isomorphism for $* < 8n-1$ and a surjection for $* = 8n-1$, and that $\mathrm{H}_*(M;\F_p) \cong 0$ for $* > 8n-1$.
The generator $x \in \pi_{4n-1} \left(\sOf \right)$ is the image of some class in $\pi_{4n-1} M$, which by abuse of notation we also denote by $x$.  We additionally abuse notation by using $J$ to denote the composite map $$M \longrightarrow \Sigma^{\infty} \Oo \mathrm{\langle 4n-1 \rangle} \stackrel{J}{\longrightarrow} \mathbb{S}.$$
\end{dfn}

\begin{lem} \label{lem:toda-main}
Let $h$ denote the canonical nullhomotopy of $2J(x)^2$ that arises from the fact that $J(x) \in \pi_{4n-1} \bS$ is an element in the odd degree homotopy of the $\E_{\infty}$-ring $\bS$.  Let $g$ denote the canonical homotopy $J(xJ(x))\simeq J(x)^2$ that arises from $J$ being a map of right $\bS$-modules, and let $f$ denote a completely arbitrary nullhomotopy of $2xJ(x)$.  Then we may form the following diagram,
$$
\begin{tikzcd}[column sep = huge]
\mathbb{S}^{8n-2} \arrow{r}{2}
\arrow[rrdd, bend right=20,""{name=D}] \arrow[rrdd, bend right=20,swap,"0"] \arrow[rr,bend left = 30,""{name=U},"0"]
\arrow[rrdd, bend right=20,""{name=D}] \arrow[rrdd, bend right=20,swap,"0"] \arrow[rr,bend left = 30,""{name=U},"0"]
& \mathbb{S}^{8n-2} \arrow[Leftrightarrow, from=D, "h"] \arrow[Leftrightarrow, from=U, "f"] \arrow{r}{xJ(x)} \arrow[rdd,""{name=MU},"J(x)^2" description] \arrow[rdd,swap, ""{name=MD},"J(x)^2" description]
& M \arrow{dd}{J}  \arrow[bend left=20,Leftrightarrow, from=MU, swap, "g"] \\ \\
& & \mathbb{S}, 
\end{tikzcd}
$$
which composes to give a Toda bracket $w \in \pi_{8n-1} (\bS)$.
The kernel of the map
$$\pi_{8n-1} (\mathbb{S}) \to \pi_{8n-1} (\MO \langle 4n \rangle)$$
is generated by $\J_{8n-1}$ and $w$. 
\end{lem}

\begin{proof}
    Since $M$ is an $(8n-1)$-skeleton, $2xJ(x)$ is trivial not just in $\pi_{8n-2} \left(\sOf \right)$, but also in $\pi_{8n-2} (M)$.  A nullhomotopy $f$ of $2x J(x)$ inside of $\pi_{8n-2} (M)$ in particular induces such a nullhomotopy in $\pi_{8n-2} \left( \sOf \right)$.  Also, the map $M \to \sOf$ is a map of right $\mathbb{S}$-modules (since it is a map of spectra), and so the homotopy $g$ from this lemma composes with the inclusion of $M$ to give the homotopy $g$ of Construction \ref{cnstr:toda-main-preskeleton}.
\end{proof}


%%% Local Variables:
%%% mode: latex
%%% eval: (visual-line-mode t)
%%% eval: (auto-fill-mode 0)
%%% End:
